\begin{Exercise}[label={ex:interessi}]
  Si progetti in C un programma che legga un float, rappresentante un ammontare di euro; di seguito il programma deve leggere un tasso d'interesse (in percentuale), ed un numero di anni. Il programma deve stampare, in uscita, per ogni anno, come l'ammontare cresce con gli interessi. Si ricordi che l'interesse si calcola con la seguente formula:
  $$
    C_{fin} = C_{in} * \left(1+\frac{r}{100}\right)^N
  $$
  Dove $C_{fin}$ è il capitale finale, $C_{in}$ è quello iniziale, $r$ è l'interesse e $N$ reppresenta il numero di anni in cui si applicano gli interessi.
\end{Exercise}

\begin{Answer}[ref={ex:interessi}]
  \lstinputlisting[style=cstyle]{src/interessi.c}
\end{Answer}
